\begin{textsample}{POS Dim 6 – mg – Score 7.00 – mg/mg\_ef\_ciencias\_da\_natureza\_intro.txt}  \label{ex:f6_pos_248}
11.1 Apresentação da Área : Ciências da Natureza A ciência está diretamente vinculada à história da humanidade.
Desde o princípio da vida humana existem práticas que revelam as inquietudes do homem na busca da solução dos \textbf{problemas} do dia a dia.
Estudar Ciências, na escola, possibilita, muitas vezes, a percepção do mundo não como verdade absoluta, mas, sim como processo de investigação, de como as relações se estabelecem, considerando que conhecer é poder olhar sobre a natureza e a vida compreendendo os processos físicos, químicos, biológicos e sociais no avanço tecnológico da humanidade.
A área de Ciências da Natureza aborda o conhecimento científico nos aspectos físicos, químicos e biológicos, por meio da investigação da natureza para interpretar de forma crítica e analítica os fenômenos naturais observados, \textbf{resultantes} das relações históricas, sociais e econômicas, visando à formação de sujeitos que atuem como agentes questionadores e transformadores, conscientes de sua responsabilidade frente aos fenômenos naturais.
Em Ciências da Natureza, os procedimentos que correspondem aos modos de buscar e organizar conhecimentos são bastante variados : a observação, a experimentação, a comparação, a elaboração de hipóteses e suposições, o debate oral sobre hipóteses, o estabelecimento de relações entre fatos ou fenômenos e ideias por meio da leitura e escrita de textos informativos, a elaboração de roteiros de pesquisa bibliográfica e questões para enquete, a busca de informações em fontes variadas, a organização de informações por meio de desenhos, tabelas, gráficos, esquemas e textos, o confronto entre suposições entre elas e os dados obtidos por investigação, a elaboração de perguntas e \textbf{problemas}, a proposição para a solução de \textbf{problemas}.
Pensando assim, o ensino de Ciências deve promover situações nas quais os estudantes possam : Definição de \textbf{Problemas} - Observar o mundo a sua volta e fazer perguntas. - Analisar demandas, delinear \textbf{problemas} e planejar investigações. - Propor hipóteses.
Levantamento, análise e representação - Planejar e realizar atividades de campo ( experimentos, observações, leituras, visitas, ambientes virtuais etc. ). - Desenvolver e utilizar ferramentas, inclusive \textbf{digitais}, para coleta, análise e representação de dados ( imagens, esquemas, tabelas, gráficos, quadros, diagramas, mapas, modelos, representações de sistemas, fluxogramas, mapas conceituais, simulações, \textbf{aplicativos} etc. ). - Avaliar informação ( validade, coerência e adequação ao \textbf{problema} formulado ). - Elaborar explicações e/ou modelos. - Associar explicações e/ou modelos à evolução histórica dos conhecimentos científicos envolvidos. - Selecionar e construir \textbf{argumentos} com base em evidências, modelos e/ou conhecimentos científicos. - Aprimorar seus saberes e incorporar, gradualmente, e de modo significativo, o conhecimento científico. - Desenvolver soluções para \textbf{problemas} cotidianos usando diferentes ferramentas, inclusive \textbf{digitais}.
Comunicação - Organizar e/ou extrapolar conclusões. - Relatar informações de forma oral, escrita ou multimodal. - Apresentar, de forma sistemática, dados e resultados de investigações. - Participar de discussões de caráter científico com colegas, professores, familiares e comunidade em geral. - Considerar contra-argumentos para rever processos investigativos e conclusões.
Intervenção - Implementar soluções e avaliar sua eficácia para resolver \textbf{problemas} cotidianos. - Desenvolver ações de intervenção para melhorar a qualidade de vida individual, coletiva e socioambiental. \textbf{Figura} 8 - Quadro de Situações \textbf{Problema} Desse modo, por meio de Unidades Temáticas, o processo de ensino-aprendizagem na área de Ciências da Natureza pode ser desenvolvido dentro de contextos sociais e culturais, que potencializam o desenvolvimento cognitivo dos estudantes, relacionando suas experiências, faixa etária, identidades culturais sociais e os diferentes significados que a ciência pode elucidar no desenvolvimento integral do ser humano.
Nesse processo, deve -se enfatizar as relações no âmbito da vida, do Universo, do ambiente e dos equipamentos tecnológicos que poderão melhorar e situar o estudante em seu mundo.
Esse desenvolvimento é concretizado quando o interesse e a curiosidade dos estudantes pela natureza, pela \textbf{tecnologia}, pelo local onde vivem e os seres que nele habitam, ou seja, o fantástico universo em que vivemos, conhecido em sua maioria, pelos meios de comunicação, despertam a curiosidade dos estudantes, favorecendo o envolvimento e o clima de interação que precisa haver para o sucesso das atividades.
Em Ciências da Natureza, o desenvolvimento de atitudes e valores envolve muitos aspectos da vida social, da cultura, do sistema produtivo e das relações entre o ser humano e a natureza.

% matched lemmas: aplicativo, argumento, digital, figurar, problema, resultante, tecnologia
\end{textsample}
